\subsection{Ambientes del ciclo de vida y promoción de componentes}

Los cinco ambientes del RT-04.01 cumplen propósitos distintos.

\begin{itemize}
  \item En \textbf{Desarrollo}, el equipo construye los módulos Laravel y las aplicaciones Kotlin, los contratos y las migraciones con datos sintéticos o anonimizados; las pruebas unitarias no necesitan conectarse al ERP real.
  \item En \textbf{QA}, se despliegan M1--M12, el verificador local y los adaptadores simulados para ensayar integración, regresión, idempotencia y fallas de enlace con datos controlados.
  \item En \textbf{Preproducción}, el mismo artefacto y las mismas versiones de contratos se ensayan con topología equivalente a producción y volumen representativo de septiembre; se prueban los 186 terminales de bodega, el relevo de turnos, la emisión de guía por el ERP y la reconciliación. Toda diferencia de escala o configuración se declara antes de aceptar el ensayo, conforme a RT-04.02.
  \item En \textbf{Producción}, los módulos sirven datos reales con acceso restringido y auditado; los desarrolladores no corrigen transacciones directamente.
  \begin{samepage}
  \item El ambiente de \textbf{Recuperación ante desastres} conserva la versión compatible de los componentes y permite ensayar conmutación y retorno al menos dos veces al año, sin presumir que una réplica remota contiene eventos aún no enviados desde un sitio aislado.
  \par\end{samepage}
\end{itemize}

Los cinco deben estar operativos en el hito H3; su emplazamiento, capacidad y segregación de redes se demuestran en 4.2, no mediante el simple nombre del ambiente.

La promoción conserva la identidad del artefacto PHP, el archivo de versiones de dependencias y la versión de esquema desde QA hasta producción.

\begin{itemize}
  \item La revisión por pares, pruebas de contrato OpenAPI/AsyncAPI, PHPUnit, PHPStan/Larastan, la auditoría de dependencias, análisis de secretos e imágenes bloquean un despliegue defectuoso (RT-04.03--04.05; Bases Técnicas Transversales, cap. 4, p. 10).
  \item Cada cambio registra requisito, commit, prueba, aprobación y despliegue; una migración incompatible con eventos sin confirmar se detiene o conserva un lector de la versión precedente.
  \item La recuperación ante desastres verifica además la lectura de datos y mensajes generados durante la contingencia antes del retorno al primario.
\end{itemize}
